% !PRINTED 756
\weberpage{756}
%<<CONT
that $C_0 F(0) + C_1 F(1) + \cdots + C_m F(m)$ is an integer not divisible by $p$, and therefore also not vanishing, and with this 1) is proved.

If we can now show further that $Q(x)$ can be made arbitrarily small for every positive $x$ by increasing $p$, it follows that for sufficiently large $p$ the expression 2) is certainly less than 1, and our proof is complete.

But if we go back to the expression (10) for $Q(x)$ and denote by $\gamma_n, \gamma_{n-1}, \ldots, \gamma_0$ the absolute values of the coefficients $c_n, c_{n-1}, \ldots, c_0$ of $f(x)$, then, since the $q_n, q_{n-1}, \ldots$ are less than 1 in absolute value, we see that for every positive $x$ the absolute value of $Q(x)$ is less than

\weq{17}{\psi(x) = \gamma_n x^n + \gamma_{n-1}x^{n-1}+\cdots+\gamma_0.}

But the coefficients $c_n, c_{n-1}, \ldots, c_0$ of the function $f(x)$ in (15) differ from the $\gamma_n, \gamma_{n-1}, \ldots, \gamma_0$ only in sign. If we therefore replace $x$ by $-x$, and so form the function
\[
f(-x) = \frac{x^{p-1}(1+x)^p(2+x)^p\ldots(m+x)^p}{\Pi(p-1)},
\]
then this function has the same coefficients as $f(x)$, but all with positive signs. It is therefore none other than the function $\psi(x)$.

If we then also set
\[
X = x(1+x)(2+x)\ldots(m+x),
\]
we obtain

\weq{18}{\psi(x) = \frac{X}{x}\cdot\frac{X^{p-1}}{\Pi(p-1)},}

and this, as was shown at the beginning of this section, approaches the limit \gesperrt{zero} as $p$ increases without bound.

\gesperrt{Thus $e$ is a transcendental number.}

\begin{center}
§. 206.

\gesperrt{Transcendence of the number $\pi$.}
\end{center}

By the same means the transcendence of the number $\pi$ can now also be proved. As the definition of this number we take the fact that it is the least positive number which satisfies the equation

\weq{1}{e^{i\pi} = -1}

when $e^x$ is defined by the formula \S.\,205, (2).
